\begin{abstract_cn}

	适当的数学模型可以准确地描述特定流行病的动态特性。
	本文旨在探索流行病学研究领域的新方法，将流行病学调查、微分方程模型和神经网络相结合，以更准确地预测和监测传染病的传播趋势。

	本文建立了一个具有年龄结构和疫苗接种的 SEIRV 仓室模型和 LSTM 模型，并讨论了流行病学调查方法、微分方程模型以及神经网络在解决流行病动力学问题方面的应用。
	通过实际案例研究，验证了这种综合方法的有效性。

	研究结果表明，将微分方程模型和神经网络相结合的方法在传染病建模中具有巨大潜力。
	这种综合性的方法能够更准确地捕捉传播动态并预测未来趋势，为应对传染病流行提供了有力的工具。
\end{abstract_cn}

\begin{keywords_cn}
	传染病模型；
	流行病学调查；
	微分方程；
	神经网络；
	COVID-19
\end{keywords_cn}